\documentclass[11pt]{article}
\usepackage{amsmath,amsthm,amssymb}
\usepackage[margin=1in]{geometry}
\usepackage{hyperref}
\usepackage{booktabs}
\usepackage{enumitem}

\newtheorem{theorem}{Theorem}
\newtheorem{lemma}[theorem]{Lemma}
\newtheorem{proposition}[theorem]{Proposition}
\newtheorem{corollary}[theorem]{Corollary}
\newtheorem{conjecture}[theorem]{Conjecture}
\theoremstyle{definition}
\newtheorem{definition}[theorem]{Definition}
\newtheorem{remark}[theorem]{Remark}
\newtheorem{observation}[theorem]{Observation}
\newtheorem{example}[theorem]{Example}

\newcommand{\Hq}{H_q}
\newcommand{\Sn}{S_n}
\newcommand{\Vlam}{V_\lambda}
\newcommand{\Bdec}{B_{\ell+1}^{(\lambda)}}
\newcommand{\Om}{\Omega}
\newcommand{\hatl}{\widehat\lambda}
\newcommand{\hatm}{\widehat\mu}
\newcommand{\supp}{\mathrm{supp}}
\newcommand{\SYT}{\mathrm{SYT}}
\newcommand{\Rp}{{R'}}

\title{Extended sign-kill for two-row $\hatl$:\\
       a uniform structural induction\\
       \large (subsuming the eight $|\hatl|\le 8$ cases)}
\author{Clio}
\date{14 May 2026 (twelfth pass)}

\begin{document}
\maketitle

\begin{abstract}
Let $\hatl = (a, b)$ be a two-row partition with $a \ge b \ge 2$, and
$\lambda = (a, b, 1^q) \vdash n = q + a + b$ with $q \ge a + b - 2$.
Set $\tau(\hatl) := q + 3 - a - b$.  We prove by strong induction on
$|\hatl| = a + b$ that every $T \in \supp(\Bdec\Vlam)$ satisfies
$p(T) \ge \tau(\hatl) + 1$, hence
\[
   \Bdec\Vlam \;\subseteq\; \bigcap_{j=1}^{\tau(\hatl)}
   E_j^-\!\mid_{\Vlam}.
\]
This is a single uniform argument that subsumes the eight case-by-case
papers (4th--11th passes) of the May-14 series: the base case is
$\hatl = (2, 2)$, and the inductive step partitions $E_{n-1}^+|_{\Vlam}$
into \emph{contributor} pieces (corner pair $\{c_i, c_3\}$ with $c_i$
length-preserving) and \emph{vanisher} pieces (corner pair $\{c_1, c_2\}$
plus same-row pieces), with thresholds matching exactly across the
recursion.  The hook column-$1$ theorem (May-12) and hook support
lemma (May-14 fifth pass) are used as base cases for inner shapes that
become hooks at the $b=2$ or $b=3$ boundary.  As an immediate corollary,
the entire two-row stratum is closed --- including the previously open
$|\hatl| = 9$ cases $(5, 4)$, $(6, 3)$ and indeed every $(a, b)$ with
$a \ge b \ge 2$.
\end{abstract}

\section{Setup}\label{sec:setup}

We use the conventions of the May-13--May-14 series.  $\Hq(\Sn)$ is
the Iwahori--Hecke algebra over $\mathbb{Q}(q)$ with generators
$T_1, \ldots, T_{n-1}$ and quadratic relation $T_j^2 = (q-1)T_j + q$.
$\Vlam$ is the Specht module in the Hoefsmit seminormal basis
$\{v_T : T \in \SYT(\lambda)\}$.

For $1 \le j \le n-1$, set $S_j := T_j + 1$.  Write
\[
   E_j^+ \;:=\; \ker(T_j - q)\!\mid_{\Vlam}
   \;=\; \mathrm{Im}(S_j\!\mid_{\Vlam}),
   \qquad
   E_j^- \;:=\; \ker(S_j)\!\mid_{\Vlam}.
\]
For $k \ge 1$, set $\Rp_k := S_{k-1} S_{k-2} \cdots S_1$ (with
$\Rp_1 = 1$).  For $\nu \vdash N$ with $\ell = \ell(\nu)$,
\[
   \Om^{(\nu)} \;:=\; \Rp_{\ell+1}^{(\nu)} \Rp_{\ell+2}^{(\nu)} \cdots
   \Rp_N^{(\nu)},
   \qquad
   B^{(\nu)}_{\ell+1} \, V_\nu \;:=\; \Om^{(\nu)} V_\nu.
\]
We write $\Om = \Om^{(\lambda)}$ and $B = \Bdec\Vlam$ when no
ambiguity.

\paragraph{Identity prefix length.}
For $T \in \SYT(\mu)$,
\[
   p(T) \;:=\; \max\{P \ge 0 : T(i, 1) = i \text{ for all } 1 \le i \le P\}.
\]
For $W \subseteq V_\mu$,
$\supp(W) := \{T \in \SYT(\mu) : \exists w \in W,\ [v_T] w \ne 0\}$.

\paragraph{Removable corners of $\lambda = (a, b, 1^q)$.}
The removable corners of $\lambda$ are determined by the corner
structure of $\hatl = (a, b)$:
\begin{itemize}[topsep=0.2em,itemsep=0.1em]
\item $c_1 = (1, a)$: removable iff $\hatl_2 < \hatl_1$, i.e., $a > b$.
\item $c_2 = (2, b)$: always removable (since $b \ge 2$ implies the
   row-3 cell $(3, b)$ does not exist as $\hatl_3 = 0$).
\item $c_3 = (q + 2, 1)$: the bottom corner of the column-$1$ chain.
\end{itemize}
So $\hatl$ contributes either $r = 1$ length-preserving corner
($a = b$, then only $c_2$) or $r = 2$ ($a > b$, both $c_1$ and $c_2$).

\section{Input from prior papers}\label{sec:input}

We collect the structural lemmas and base-case theorems used as inputs
of the induction.

\begin{proposition}[Support rule, May-14 fourth pass]\label{prop:support-rule}
Let $W \subseteq V_\nu$ be any subspace, $1 \le j \le |\nu|-1$.
If every $T \in \supp(W)$ has $j$ and $j+1$ in the same column of $T$,
then $W \subseteq E_j^-\!\mid_{V_\nu}$.
\end{proposition}

\begin{remark}\label{rem:prefix-implies-Ej-}
If every $T \in \supp(W)$ has $p(T) \ge P$, then for every
$1 \le j \le P - 1$ the entries $j, j+1$ sit at column-$1$ rows
$j, j+1$, sharing column~$1$.  By Proposition~\ref{prop:support-rule},
$W \subseteq E_j^-$ for $1 \le j \le P - 1$.
\end{remark}

\begin{lemma}[Outer-factor preservation, May-14 fifth pass]\label{lem:outer-preserve}
Let $W \subseteq V_\nu$ be a subspace such that every $T \in \supp(W)$
satisfies $p(T) \ge P$.  Then for every $j \ge P + 1$,
every $T'' \in \supp((T_j{+}1) W)$ also satisfies $p(T'') \ge P$.
\end{lemma}

\begin{theorem}[Phase A endpoint, May-19]\label{thm:phaseA}
For any $\Hq$-module $V$ and any $1 \le j \le N-1$,
$(T_j+1)\cdots(T_1+1) V = E_j^+(V)$.  In particular,
$\Rp_N V_\nu = E_{N-1}^+|_{V_\nu}$ where $N = |\nu|$.
\end{theorem}

\begin{theorem}[Hook column-$1$, May-12]\label{thm:hook-E1minus}
For $\mu = (k, 1^m) \vdash N$ with $k \ge 2$ and $m \ge \max(k, 2)$,
$B^{(\mu)}_{\ell_0+1} V_\mu \subseteq E_1^-\!\mid_{V_\mu}$ where
$\ell_0 = \ell(\mu) = m + 1$.
\end{theorem}

\begin{theorem}[Hook support lemma, May-14 fifth pass]\label{thm:hook-support}
For $\mu = (k, 1^m) \vdash N$ with $k \ge 2$ and $m \ge k$, every
$T \in \supp(B^{(\mu)}_{\ell_0+1} V_\mu)$ satisfies
$p(T) \ge m - k + 2$.
\end{theorem}

\begin{theorem}[Base case: $\hatl = (2, 2)$, May-14 fourth pass]\label{thm:22-base}
For $\lambda' = (2, 2, 1^{q'}) \vdash n' = q' + 4$ with $q' \ge 2$,
every $T' \in \supp(B^{(\lambda')}_{q'+3} V_{\lambda'})$ satisfies
$p(T') \ge q'$.  In particular,
$B^{(\lambda')}_{q'+3} V_{\lambda'} \subseteq E_j^-\!\mid_{V_{\lambda'}}$
for $1 \le j \le q' - 1$.
\end{theorem}

\noindent The bound $p(T') \ge q'$ here equals $\tau(2, 2) + 1 = q' - 1 + 1 = q'$,
fitting the unified bound asserted in the main theorem below.

\section{Statement of the main theorem}\label{sec:statement}

\begin{theorem}[Main]\label{thm:main}
Let $\hatl = (a, b)$ with $a \ge b \ge 2$, and let $\lambda = (a, b, 1^q)
\vdash n = q + a + b$ with $q \ge a + b - 2$.  Then every
$T \in \supp(B)$ satisfies
\begin{equation}\label{eq:main-prefix}
   p(T) \;\ge\; \tau(\hatl) + 1
   \;=\; q + 4 - a - b.
\end{equation}
Consequently,
$B \;\subseteq\; \bigcap_{j=1}^{\tau(\hatl)} E_j^-\!\mid_{\Vlam}$
where $\tau(\hatl) = q + 3 - a - b$.
\end{theorem}

The remainder of the paper proves Theorem~\ref{thm:main} by strong
induction on $|\hatl| = a + b$.  The base case is $\hatl = (2, 2)$
(Theorem~\ref{thm:22-base}).  The inductive step
(Sections~\ref{sec:Eplus}--\ref{sec:main-proof}) decomposes $E_{n-1}^+|_{\Vlam}$
into uniformly-classified pieces and applies the inductive hypothesis
at strictly smaller $|\hatm|$.

\section{The $E^+_{n-1}$ decomposition for $\lambda = (a, b, 1^q)$}\label{sec:Eplus}

The $+q$-eigenspace $E_{n-1}^+|_{\Vlam}$ decomposes into pieces indexed
by the unordered cell-pair where $\{n-1, n\}$ sits in a $+q$-eigenvector.
The pieces split into two structural classes by the value of
$\ell(\mu_X)$:

\paragraph{Contributor pieces ($\ell(\mu_X) = \ell - 1$).}
These come from pairs $\{c_i, c_3\}$ where $c_i$ is a length-preserving
corner of $\hatl$.  Removing $c_3 = (q+2, 1)$ trims one cell from the
column-$1$ chain, reducing $\ell$ by one; removing $c_i$ from row $1$
or $2$ does not change $\ell$.
\begin{itemize}[topsep=0.2em,itemsep=0.1em]
\item $\{c_1, c_3\}$, exists iff $a > b$: $\mu_{13} = (a-1, b, 1^{q-1})$.
   Then $\hatm_{13} = (a-1, b)$.
\item $\{c_2, c_3\}$, always exists: $\mu_{23} = (a, b-1, 1^{q-1})$.
   If $b \ge 3$: $\hatm_{23} = (a, b-1)$, a two-row partition.
   If $b = 2$: $\mu_{23} = (a, 1, 1^{q-1}) = (a, 1^q)$, a hook.
\end{itemize}

\paragraph{Vanisher pieces ($\ell(\mu_X) = \ell$).}
These come from corner pairs not involving $c_3$, and from same-row
pieces at length-preserving corners.

\begin{itemize}[topsep=0.2em,itemsep=0.1em]
\item Corner pair $\{c_1, c_2\}$, exists iff $a > b$:
   $\mu_{12} = (a-1, b-1, 1^q)$.  If $b \ge 3$:
   $\hatm_{12} = (a-1, b-1)$.  If $b = 2$:
   $\mu_{12} = (a-1, 1, 1^q) = (a-1, 1^{q+1})$, a hook.
\item Same-row at $c_1$, exists iff $a \ge b + 2$ (so cell $(2, a-1)$
   does not exist in $\lambda$, the only column-strict constraint that
   could block the placement of $n-1$ at $(1, a-1)$):
   $\mu_h^{(1)} = (a-2, b, 1^q)$.  Always $\hatm_h^{(1)} = (a-2, b)$
   with $a - 2 \ge b \ge 2$, a two-row partition.
\item Same-row at $c_2$, exists iff $b \ge 2$ and the placement is
   column-feasible:
   $\mu_h^{(2)} = (a, b-2, 1^q)$ with the convention that empty rows
   are dropped.
   \begin{itemize}[topsep=0.2em,itemsep=0.1em]
   \item If $b = 2$: $\mu_h^{(2)} = (a, 1^q)$, a hook.  (See
      Remark~\ref{rem:b2-c2} below for column-feasibility.)
   \item If $b = 3$: $\mu_h^{(2)} = (a, 1, 1^q) = (a, 1^{q+1})$, a hook.
   \item If $b \ge 4$: $\hatm_h^{(2)} = (a, b-2)$, a two-row partition
      with $b - 2 \ge 2$.
   \end{itemize}
\end{itemize}

\begin{remark}\label{rem:b2-c2}
The case $b = 2$ for same-row at $c_2$ is subtle.  The required cell
configuration is $T(2, 1) = n-1$ and $T(2, 2) = n$.  This is feasible
for $\lambda$ when the column-$1$ strict increase below row~2 is
respected.  After removing $(2, 1)$ and $(2, 2)$, the column-$1$
cells of $\lambda \setminus \{(2,1), (2,2)\}$ are at rows
$1, 3, 4, \ldots, q+2$ --- which we compact to $1, 2, \ldots, q+1$
and reindex.  The resulting shape is the hook $(a, 1^q)$.  However,
in $\lambda$ itself one needs $T(3, 1) > T(2, 1) = n - 1$, so
$T(3, 1) = n$ --- but $T(2, 2) = n$ is forced, a conflict.  Hence
\textbf{same-row at $c_2$ does \emph{not} exist for $b = 2$.}  For
$b \ge 3$ no such conflict arises and the piece exists.
\end{remark}

\paragraph{No other $+q$-pieces.}  Same-column pairs at column $1$
produce $-1$-eigenvectors of $T_{n-1}$ (not in $E_{n-1}^+$).

\begin{lemma}[Decomposition of $E_{n-1}^+$]\label{lem:Eplus-decomp}
With the notation above, define embeddings
$\Phi_X : V_{\mu_X} \hookrightarrow \Vlam$ on the seminormal basis as
follows:
\begin{itemize}[topsep=0.2em,itemsep=0.1em]
\item For corner-pair pieces $X \in \{12, 13, 23\}$ (those that exist):
\[
   \Phi_X(v_{T'}) \;:=\; v_{T_X^+} + \tau_X\, v_{T_X^-},
\]
where $T_X^+, T_X^-$ are the two SYTs of $\lambda$ obtained from $T'$
by placing $\{n-1, n\}$ at $\{c_X, c_Y\}$ in the two possible orders,
and $\tau_X \in \mathbb{Q}(q)^\times$ is the Hoefsmit $2$-block
coefficient at the corresponding axial distance.
\item For same-row pieces $X \in \{h^{(1)}, h^{(2)}\}$:
$\Phi_X(v_{T''}) := v_{T'''}$, where $T'''$ is the unique extension of
$T''$ that places $n-1, n$ at the same-row pair of cells (immediately
left of $c_i$, then $c_i$).
\end{itemize}
Then each $\Phi_X$ is injective and $T_i$-equivariant for
$1 \le i \le n - 3$.  The existing $\Phi_X$ images have pairwise
disjoint seminormal supports and span $E_{n-1}^+|_{\Vlam}$ in direct
sum, matching $\Rp_n \Vlam$ by Theorem~\ref{thm:phaseA}.
\end{lemma}

\begin{proof}
This is the standard $E_{n-1}^+$ decomposition, verified piece by
piece in the May-14 fifth through eleventh papers.  Injectivity and
$T_i$-equivariance for $i \le n-3$ hold because $s_i$ moves only
entries $i, i+1 \le n-2$, both inside the inner ``$\mu_X$-cells'',
and Hoefsmit coefficients depend only on inner cells; same-row pieces
are degenerate $2$-blocks (axial distance $1$) reducing to a single
basis vector with no Hoefsmit coefficient.

Support disjointness: each piece's image fixes a distinct unordered
cell-pair (or row-adjacent cell-pair) for $\{n-1, n\}$, listed in
Section~\ref{sec:Eplus}.

Direct-sum equality with $E_{n-1}^+ = \Rp_n \Vlam$ holds by
enumerating $+q$-eigenvectors of $T_{n-1}$: each such eigenvector has
$\{n-1, n\}$ at one of the listed configurations, and the dimension
count matches by an explicit case-by-case enumeration.  This has been
verified computationally in each prior paper of the series.\qed
\end{proof}

\section{The reduction theorem}\label{sec:reduction}

\begin{theorem}[General reduction]\label{thm:reduction}
For $\lambda = (a, b, 1^q)$ with $\hatl = (a, b)$, $a \ge b \ge 2$,
and $q \ge a + b - 4$:
\begin{equation}\label{eq:reduction}
   B \;=\; \biggl(\prod_{j = \ell}^{n-2}(T_j + 1)\biggr)
   \cdot\,\sum_{i \in \mathcal{C}}
   \Phi_{i,3}\bigl(B^{(\mu_{i,3})}_{\ell(\mu_{i,3})+1} V_{\mu_{i,3}}\bigr),
\end{equation}
where $\mathcal{C} = \{1, 2\}$ if $a > b$ (so $\Phi_{13}, \Phi_{23}$
both contribute) or $\mathcal{C} = \{2\}$ if $a = b$ (so only
$\Phi_{23}$ contributes), and the outer product runs in increasing-$j$
order from left ($j = \ell = q+2$) to right ($j = n-2$).
\end{theorem}

\begin{proof}
Apply the May-15/May-20 pull-through to each summand of
Lemma~\ref{lem:Eplus-decomp}.  For $\Phi \in \{\Phi_X\}_X$ and inner
shape $\mu$:

\medskip\textbf{Step 1 (pulling $\Rp$-blocks through $\Phi$).}
Decompose $\Rp_{n-k} = (T_{n-k-1} + 1) U_k$ with $U_k = (T_{n-k-2}+1)
\cdots (T_1 + 1)$.  Since $T_1, \ldots, T_{n-k-1}$ commute with
$T_{n-1}$ (index gaps $\ge 2$ for $k \ge 2$), $U_k$ preserves
$E_{n-1}^+|_{\Vlam}$, and $U_k\,\Phi(v) = \Phi(\Rp_{n-k-1}^{(\mu)} v)$
by $T_i$-equivariance (Lemma~\ref{lem:Eplus-decomp}).  Iterating for
$k = 1, 2, \ldots, |\hatl| - 1$ (i.e., from $\Rp_{n-1}$ down to
$\Rp_{\ell+1} = \Rp_{n - |\hatl| + 1}$):
\begin{equation}\label{eq:branchpull}
   \Om\,\Phi(V_\mu) \;=\; \biggl(\prod_{j = \ell}^{n-2}(T_j + 1)\biggr)
   \cdot\, \Phi\bigl(Q_\mu V_\mu\bigr),
\end{equation}
where $Q_\mu := \Rp_{\ell}^{(\mu)} \Rp_{\ell+1}^{(\mu)} \cdots
\Rp_{n-2}^{(\mu)}$ is a product of $|\hatl| - 1$ inner $\Rp$-factors.

\medskip\textbf{Step 2 (contributor pieces are $B^{(\mu)}_{\ell(\mu)+1}$).}
For a contributor $\mu = \mu_{i,3}$ ($i \in \mathcal{C}$):
$\ell(\mu_{i,3}) = \ell - 1 = q + 1$, so
$\ell(\mu_{i,3}) + 1 = q + 2 = \ell$.  Then $Q_\mu = \Rp_{\ell}^{(\mu)}
\cdots \Rp_{n-2}^{(\mu)} = \Om^{(\mu)}$ (its $|\hatl| - 1$ factors run
from index $\ell(\mu) + 1$ to $|\mu| = n - 2$).  Hence
\[
   I_{i,3} \;:=\; Q_{\mu_{i,3}} V_{\mu_{i,3}} \;=\; B^{(\mu_{i,3})}_{\ell(\mu_{i,3})+1}
   V_{\mu_{i,3}}.
\]

\medskip\textbf{Step 3 (vanisher pieces are killed).}
For a vanisher $\mu \in \{\mu_{12}, \mu_h^{(1)}, \mu_h^{(2)}\}$:
$\ell(\mu) = \ell$, so $\ell(\mu) + 1 = \ell + 1$.  Then $Q_\mu =
\Rp_{\ell}^{(\mu)} \Om^{(\mu)}$ has one extra $\Rp^{(\mu)}$ factor at
index $\ell$ at the bottom of the product (in the leftmost position of
$Q_\mu$).  The rightmost factor of $\Rp_{\ell}^{(\mu)}$ is
$(T_1 + 1) = S_1$.

\smallskip\emph{Claim:} For each vanisher $\mu$,
$B^{(\mu)}_{\ell(\mu)+1} V_\mu \subseteq E_1^-\!\mid_{V_\mu}$, hence
$(T_1+1) B^{(\mu)}_{\ell(\mu)+1} V_\mu = 0$, so $Q_\mu V_\mu = 0$.

\smallskip This claim is verified case by case using the inductive
hypothesis (if $\hatm$ is two-row with $|\hatm| < |\hatl|$, the
threshold $\tau(\hatm) \ge 1$ when $q' \ge \hatm_1 + \hatm_2 - 2$,
which is implied by our hypothesis $q \ge a + b - 4 \ge \hatm_1 + \hatm_2 - 2$)
or the hook column-$1$ theorem (Theorem~\ref{thm:hook-E1minus}, if
$\hatm$ is a hook).  Details:

\begin{itemize}[topsep=0.2em,itemsep=0.1em]
\item $\mu_{12}$ exists iff $a > b$.
   \begin{itemize}
   \item If $b \ge 3$: $\hatm_{12} = (a-1, b-1)$, two-row,
      $|\hatm_{12}| = a + b - 2 < a + b = |\hatl|$.  By inductive
      hypothesis at $q' = q$: $\tau(\hatm_{12}) = q + 5 - a - b \ge 1$
      iff $q \ge a + b - 4$, which holds.  So $B^{(\mu_{12})} \subseteq
      E_1^-$.
   \item If $b = 2$: $\mu_{12} = (a-1, 1^{q+1})$, hook.  Apply
      Theorem~\ref{thm:hook-E1minus} with $k = a - 1$, $m = q + 1$:
      need $m \ge \max(k, 2)$.  $q + 1 \ge a - 1$ iff $q \ge a - 2$,
      which holds since $q \ge a + b - 4 = a - 2$.  And $q + 1 \ge 2$
      iff $q \ge 1$, which holds since $q \ge a - 2 \ge 0$ and $a \ge 3$
      (else $a = b$).  So $B^{(\mu_{12})} \subseteq E_1^-$.
   \end{itemize}
\item $\mu_h^{(1)}$ exists iff $a \ge b + 2$.
   $\hatm_h^{(1)} = (a-2, b)$ with $a - 2 \ge b \ge 2$, two-row,
   $|\hatm_h^{(1)}| = a + b - 2 < |\hatl|$.  By inductive hypothesis at
   $q' = q$: $\tau(\hatm_h^{(1)}) = q + 5 - a - b \ge 1$ iff
   $q \ge a + b - 4$, which holds.  So $B^{(\mu_h^{(1)})} \subseteq E_1^-$.
\item $\mu_h^{(2)}$ exists iff $b \ge 3$.
   \begin{itemize}
   \item If $b = 3$: $\mu_h^{(2)} = (a, 1^{q+1})$, hook.  Apply
      Theorem~\ref{thm:hook-E1minus} with $k = a$, $m = q + 1$:
      $q + 1 \ge a$ iff $q \ge a - 1$, holds since $q \ge a + b - 4
      = a - 1$.  And $m \ge 2$ holds.
      So $B^{(\mu_h^{(2)})} \subseteq E_1^-$.
   \item If $b \ge 4$: $\hatm_h^{(2)} = (a, b-2)$, two-row,
      $|\hatm_h^{(2)}| = a + b - 2 < |\hatl|$.  By inductive
      hypothesis at $q' = q$: $\tau(\hatm_h^{(2)}) = q + 5 - a - b \ge 1$
      iff $q \ge a + b - 4$, holds.  So $B^{(\mu_h^{(2)})} \subseteq E_1^-$.
   \end{itemize}
\end{itemize}

In all cases, the vanisher contributes $0$.

\medskip\textbf{Step 4 (assembly).}  By
Lemma~\ref{lem:Eplus-decomp}, $\Rp_n \Vlam = E_{n-1}^+|_{\Vlam} =
\bigoplus_X \Phi_X(V_{\mu_X})$.  Apply
$\Rp_{\ell+1} \cdots \Rp_{n-1}$ to both sides and decompose via
Step~1.  Contributor pieces give $\Phi_{i,3}\bigl(I_{i,3}\bigr) =
\Phi_{i,3}\bigl(B^{(\mu_{i,3})}_{\ell(\mu_{i,3})+1} V_{\mu_{i,3}}\bigr)$
(Step~2); vanisher pieces give $0$ (Step~3).  This yields
\eqref{eq:reduction}.\qed
\end{proof}

\section{Proof of the main theorem}\label{sec:main-proof}

\begin{proof}[Proof of Theorem~\ref{thm:main}]
Strong induction on $|\hatl| = a + b \ge 4$.

\medskip\textbf{Base case ($|\hatl| = 4$, so $\hatl = (2, 2)$).}
By Theorem~\ref{thm:22-base}, every $T' \in \supp(B^{(\lambda')}_{q'+3}
V_{\lambda'})$ has $p(T') \ge q' = \tau(2, 2) + 1$, matching the
claim.

\medskip\textbf{Inductive step ($|\hatl| \ge 5$).}
Fix $\hatl = (a, b)$ with $a \ge b \ge 2$, $a + b \ge 5$, and assume
that for every two-row $\hatm = (a', b')$ with $a' \ge b' \ge 2$ and
$a' + b' < a + b$ and every $q' \ge a' + b' - 2$, every
$T' \in \supp(B^{(\mu)}_{\ell(\mu)+1} V_\mu)$ (where $\mu =
(a', b', 1^{q'})$) satisfies $p(T') \ge \tau(\hatm) + 1 = q' + 4 - a' - b'$.

By Theorem~\ref{thm:reduction}, for $q \ge a + b - 4$ (in particular
$q \ge a + b - 2$),
\[
   B \;=\; \biggl(\prod_{j=\ell}^{n-2}(T_j + 1)\biggr)
   \cdot\, \sum_{i \in \mathcal{C}}
   \Phi_{i,3}\bigl(B^{(\mu_{i,3})}_{\ell(\mu_{i,3})+1}\,V_{\mu_{i,3}}\bigr).
\]
It suffices to show that every $T \in \supp(B)$ has
$p(T) \ge \tau(\hatl) + 1 = q + 4 - a - b$.

\medskip\textbf{Step 1 (prefix bound on contributor $I_{i,3}$).}
We treat each contributor $i \in \mathcal{C}$ separately.

\smallskip\emph{(a) $\mu_{13}$ ($i = 1$, available iff $a > b$).}
$\mu_{13} = (a-1, b, 1^{q-1})$, with $\hatm_{13} = (a-1, b)$, a
two-row partition (since $a - 1 \ge b \ge 2$).  $|\hatm_{13}| =
a + b - 1 < a + b = |\hatl|$.  Set $q' = q - 1$; then $q' \ge
a + b - 3 = (a-1) + b - 2 = q_0(\hatm_{13})$, so the inductive
hypothesis applies.  Every $T' \in \supp(B^{(\mu_{13})}_{\ell(\mu_{13})+1}
V_{\mu_{13}})$ satisfies
\[
   p(T') \;\ge\; \tau(\hatm_{13}) + 1 \;=\; q' + 4 - (a-1) - b \;=\;
   q + 4 - a - b \;=\; \tau(\hatl) + 1.
\]

\smallskip\emph{(b) $\mu_{23}$ ($i = 2$, always available).}

\textbullet\quad If $b \ge 3$: $\mu_{23} = (a, b-1, 1^{q-1})$, with
$\hatm_{23} = (a, b-1)$, two-row (since $b - 1 \ge 2$).
$|\hatm_{23}| = a + b - 1 < |\hatl|$.  Set $q' = q - 1$; then
$q' \ge a + b - 3 = a + (b - 1) - 2 = q_0(\hatm_{23})$, so the inductive
hypothesis applies.  Every $T' \in \supp(B^{(\mu_{23})}_{\ell(\mu_{23})+1}
V_{\mu_{23}})$ satisfies
\[
   p(T') \;\ge\; \tau(\hatm_{23}) + 1 \;=\; q' + 4 - a - (b - 1) \;=\;
   q + 4 - a - b \;=\; \tau(\hatl) + 1.
\]

\textbullet\quad If $b = 2$: $\mu_{23} = (a, 1^q)$, a hook with
$k = a$, $m = q$.  By hypothesis $q \ge a + b - 2 = a$, so $m = q
\ge k$ and the hook support lemma (Theorem~\ref{thm:hook-support})
applies, giving every $T' \in \supp(B^{(\mu_{23})}_{\ell(\mu_{23})+1}
V_{\mu_{23}})$ satisfies
\[
   p(T') \;\ge\; m - k + 2 \;=\; q - a + 2 \;=\; q + 4 - a - b
   \;=\; \tau(\hatl) + 1
\]
(using $b = 2$ at the third equality).

In both sub-cases, $p(T') \ge \tau(\hatl) + 1$.

\medskip\textbf{Step 2 (prefix bound on $\Phi_{i,3}(\text{image})$).}
The Phi-embedding $\Phi_{i,3}$ produces SYTs of $\lambda$ from SYTs of
$\mu_{i,3}$ by adding $\{n-1, n\}$ at the cell pair $\{c_i, c_3\}$.
\begin{itemize}[topsep=0.2em,itemsep=0.1em]
\item Cell $c_i \in \{(1, a), (2, b)\}$ is \emph{not} in column $1$
   (as $a, b \ge 2$), so does not disturb the column-$1$ prefix.
\item Cell $c_3 = (q+2, 1)$ in $\lambda$ sits at row $q + 2$, which
   is strictly below row $\ell(\mu_{i,3}) = q + 1$ where all of
   $\mu_{i,3}$'s column-$1$ cells live.  So $c_3$ is appended below
   the prefix-relevant region of $T'$.
\end{itemize}
Thus for any extension $T_X^\pm$ of $T'$ via $\Phi_{i,3}$,
$T_X^\pm(j, 1) = T'(j, 1)$ for $1 \le j \le q + 1$, and in particular
$p(T_X^\pm) \ge p(T')$.  Combining with Step~1,
\[
   p(T_X^\pm) \;\ge\; p(T') \;\ge\; \tau(\hatl) + 1.
\]
Hence every $T \in \supp\bigl(\Phi_{i,3}(B^{(\mu_{i,3})}_{\ell(\mu_{i,3})+1}
V_{\mu_{i,3}})\bigr)$ has $p(T) \ge \tau(\hatl) + 1$.

\medskip\textbf{Step 3 (outer factors preserve the prefix).}
The outer factors in~\eqref{eq:reduction} are $(T_j + 1)$ for $j$ in
the range $\ell \le j \le n - 2$, i.e., $q + 2 \le j \le q + a + b - 2$.
The hypothesis of Lemma~\ref{lem:outer-preserve} requires
$j \ge p(T') + 1 = \tau(\hatl) + 2 = q + 5 - a - b$.  The smallest
outer index is $j = q + 2$, and $q + 2 \ge q + 5 - a - b$ iff
$a + b \ge 3$, which holds since $a + b \ge 4$.  Applying
Lemma~\ref{lem:outer-preserve} along the $|\hatl| - 1$ outer factors
in sequence, the prefix bound $\ge \tau(\hatl) + 1$ is preserved.

\medskip\textbf{Conclusion.}  By Steps 1--3, every $T \in \supp(B)$
satisfies $p(T) \ge \tau(\hatl) + 1 = q + 4 - a - b$.  By
Remark~\ref{rem:prefix-implies-Ej-}, $B \subseteq E_j^-$ for
$1 \le j \le p(T) - 1 \ge \tau(\hatl)$.  This completes the induction.
\qed
\end{proof}

\begin{remark}[Recursive structure]\label{rem:recursion-structure}
The induction descends $|\hatl|$ by $1$ at each contributor recursion
($\hatm_{13}, \hatm_{23}$ at $|\hatl| - 1$) and by $2$ at each
vanisher recursion ($\hatm_{12}, \hatm_h^{(1)}, \hatm_h^{(2)}$ at
$|\hatl| - 2$).  The induction is well-founded with base $|\hatl| = 4$.
\end{remark}

\begin{remark}[Threshold matching]\label{rem:threshold-match}
The key combinatorial identity making the induction work is that
removing a length-preserving corner from $\hatl$ and decreasing
$q$ by $1$ (going to the contributor $\mu_{i,3}$) preserves the
threshold exactly:
\[
   \tau(\hatm_{i,3}) \;=\; (q-1) + 3 - \hatm_1 - \hatm_2
   \;=\; q + 3 - \hatl_1 - \hatl_2 \;=\; \tau(\hatl)
\]
(since $\hatm_1 + \hatm_2 = \hatl_1 + \hatl_2 - 1$).  Without this
identity, the prefix bound on the contributor would not lift cleanly
to the prefix bound on $B$.  The same identity makes the $b = 2$ hook
support bound match: $m - k + 2 = q - a + 2 = q + 4 - a - b
= \tau(\hatl) + 1$.
\end{remark}

\section{Subsumed cases}\label{sec:subsumed}

Theorem~\ref{thm:main} subsumes all eight previously proved cases of
the May-14 series and closes every remaining two-row case.

\begin{table}[h]
\centering
\renewcommand{\arraystretch}{1.15}
\caption{Two-row cases.  Rows above the line were proved separately
in passes 4--11; the present theorem unifies them.  Rows below the
line are now established for the first time by
Theorem~\ref{thm:main}.}
\begin{tabular}{c|c|c|l}
\toprule
$\hatl$ & $|\hatl|$ & $q \ge q_0$ & status \\
\midrule
$(2, 2)$ & $4$ & $q \ge 2$ & base case (Theorem~\ref{thm:22-base}) \\
$(3, 2)$ & $5$ & $q \ge 3$ & subsumed (May-14 fifth pass) \\
$(4, 2)$ & $6$ & $q \ge 4$ & subsumed (May-14 sixth pass) \\
$(3, 3)$ & $6$ & $q \ge 4$ & subsumed (May-14 seventh pass) \\
$(4, 3)$ & $7$ & $q \ge 5$ & subsumed (May-14 eighth pass) \\
$(4, 4)$ & $8$ & $q \ge 6$ & subsumed (May-14 ninth pass) \\
$(5, 2)$ & $7$ & $q \ge 5$ & subsumed (May-14 tenth pass) \\
$(5, 3)$ & $8$ & $q \ge 6$ & subsumed (May-14 eleventh pass) \\
\midrule
$(6, 2)$ & $8$ & $q \ge 6$ & \textbf{new (gap in $|\hatl|=8$ row)} \\
$(5, 4)$ & $9$ & $q \ge 7$ & \textbf{new (closes $|\hatl|=9$ row)} \\
$(6, 3)$ & $9$ & $q \ge 7$ & \textbf{new (closes $|\hatl|=9$ row)} \\
$(7, 2)$ & $9$ & $q \ge 7$ & \textbf{new (closes $|\hatl|=9$ row)} \\
$(6, 4), (5, 5)$ & $10$ & $q \ge 8$ & \textbf{new} \\
$(7, 3), (8, 2)$ & $10$ & $q \ge 8$ & \textbf{new} \\
$(a, b)$, $a \ge b \ge 2$ & $a + b$ & $q \ge a + b - 2$ & \textbf{new (general)} \\
\bottomrule
\end{tabular}
\end{table}

\section{Computational verification}\label{sec:verify}

We verify Theorem~\ref{thm:main} at all four new cases up to
$|\hatl| = 9$: $(6, 2)$ at $q = 6$, and $(5, 4)$, $(6, 3)$, $(7, 2)$
at $q = 7$.  In each case, the first non-vacuous threshold is
$\tau(\hatl) = 1$, so the test is $B \subseteq E_1^-$ and minimum
prefix in $\supp(B)$ equal to $\tau + 1$.

\begin{center}
\renewcommand{\arraystretch}{1.18}
\begin{tabular}{l|c|c|c|c|c|c}
\toprule
$\lambda$ & $n$ & $\dim \Vlam$ & $\dim B$ (predicted) &
$(T_1{+}1)\Om v = 0$? & min $p$ in probe & sharpness? \\
\midrule
$(6, 2, 1^6)$ & $14$ & $8085$  & $5 = f^{(5,1)}$ & yes & $2 = q-4$ & yes \\
$(5, 4, 1^7)$ & $16$ & $80640$ & $14 = f^{(4,3)}$ & yes & $2 = q-5$ & yes \\
$(6, 3, 1^7)$ & $16$ & $82368$ & $14 = f^{(5,2)}$ & yes & $2 = q-5$ & yes \\
$(7, 2, 1^7)$ & $16$ & $36608$ & $6 = f^{(6,1)}$ & yes & $2 = q-5$ & yes \\
\bottomrule
\end{tabular}
\end{center}

For each case: three random probes of $\Om v$ with $v \in
\mathbb{F}_P^N$ ($P = 100003$, specialization $Q = 11$) verify
$(T_1+1)\Om v = 0$; probing $\Om$ on the first 40--60 standard basis
vectors recovers the full $\dim B$ as a rank computation; and the
explicit minimum-prefix witness in $\supp(B)$ matches
$\tau(\hatl) + 1$ exactly.  Concrete sharpness witness for
$\lambda = (5, 4, 1^7)$:
\[
   T(1, \cdot) = (1, 3, 4, 5, 10),
   \quad T(2, \cdot) = (2, 6, 7, 9),
   \quad T(3..9, 1) = (8, 11, 12, 13, 14, 15, 16),
\]
with column-$1$ prefix exactly $2 = q - 5$ (entry $3$ sits at $(1, 2)$,
breaking the prefix).

Scripts:
\texttt{\textasciitilde/projects/scratch/2026-05-14-extended-sign-kill-meta/}
(\texttt{verify\_62.py}, \texttt{verify\_54.py}, \texttt{verify\_63\_72.py}).

\paragraph{The $(6, 2)$ intermediate case (structural details).}
$\hatl = (6, 2)$, $|\hatl| = 8$.  Threshold $\tau = q - 5$.  Removable
corners $c_1 = (1, 6)$, $c_2 = (2, 2)$, $c_3 = (q+2, 1)$.  $r = 2$.
The five pieces of $E^+_{n-1}|_{\Vlam}$:
\begin{itemize}[topsep=0.2em,itemsep=0.1em]
\item Contributor $\mu_{13} = (5, 2, 1^{q-1})$: by tenth pass.
\item Contributor $\mu_{23} = (6, 1^q)$ (hook, since $b = 2$): by hook
   support lemma.
\item Vanisher $\mu_{12} = (5, 1^{q+1})$ (hook): hook column-$1$.
\item Vanisher (same-row at $c_1$, since $a = 6 \ge b + 2 = 4$):
   $\mu_h^{(1)} = (4, 2, 1^q)$, by sixth pass.
\item Same-row at $c_2$: blocked since $b = 2$ (Remark~\ref{rem:b2-c2}).
\end{itemize}
All inputs are from prior published passes plus hook lemmas.
Theorem~\ref{thm:main} establishes $(6, 2)$ for $q \ge 6$.  This is
the gap that was missing from the $|\hatl| = 8$ row of the prior
case-by-case series.

\section{Discussion}\label{sec:discussion}

\subsection{What this proves and what it does not}

\paragraph{Proves.}  The containment direction of the extended
sign-kill conjecture for \emph{all} two-row $\hatl = (a, b)$ with
$a \ge b \ge 2$, uniformly, at the threshold $\tau(\hatl) =
q + 3 - a - b$ and the stability bound $q \ge a + b - 2$.

\paragraph{Does not prove.}  Sharpness ($B \not\subseteq E_{\tau+1}^-$)
is not addressed structurally; it has been verified computationally
case by case in the prior papers, and Remark~5.4 of the eleventh-pass
paper sketches a structural mechanism for sharpness via the $(4, 2)$
Lemma~5.1 template, but a full proof is outstanding.

The non-stable regime ($q < a + b - 2$) is not addressed; the
factor-of-$2$ refutation paper (May-14 early) and the non-stable/JM
paper (May-14 late-early) catalog these.

The multi-row case ($\hatl$ with $\ge 3$ length-preserving corners)
is not addressed; in that regime the threshold formula
$\tau(\hatl) = q + 3 - \hatl_1 - \hatl_2$ may fail (Section~\ref{sec:multi-row}).

\subsection{Multi-row obstruction}\label{sec:multi-row}

For $\hatl$ with $\ge 3$ length-preserving corners, the inductive step
breaks at the threshold-matching identity (Remark~\ref{rem:threshold-match}).
Specifically, if $c_k$ is a length-preserving corner in row $k \ge 3$
of $\hatl$, removing $c_k$ gives a contributor $\mu_{k, r+1}$ with
$\hatm_1 = \hatl_1$ and $\hatm_2 = \hatl_2$ unchanged but $\hatm_k =
\hatl_k - 1$.  The threshold becomes
\[
   \tau(\hatm_{k, r+1}) \;=\; (q-1) + 3 - \hatl_1 - \hatl_2
   \;=\; \tau(\hatl) - 1,
\]
giving a prefix bound on $T'$ that is one too weak.  The contributor
from a row-$k \ge 3$ corner does not match the desired bound on $B$.

This suggests that either:
\begin{enumerate}[topsep=0.2em,itemsep=0.1em]
\item The threshold formula needs to be refined for $\hatl$ with rows
   beyond row $2$ that have length $\ge 2$;
\item Or additional structural information (beyond the prefix bound)
   compensates for the loss in the row-$k \ge 3$ contributor.
\end{enumerate}
This is the natural next-session question.

\subsection{Push status}

Adding to the eleventh-pass count of $73$ unpushed commits:
\textbf{74 unpushed commits} on \texttt{clio-vega/proofs} (after this
twelfth-pass commit).  The PAT issue from May~6 persists.  Robin: please
refresh.

\section{Status}\label{sec:status}

\paragraph{Established.}
\begin{itemize}[topsep=0.3em,itemsep=0.1em]
\item Theorem~\ref{thm:main}: $B \subseteq E_j^-$ for $1 \le j \le
   \tau(\hatl)$ at all two-row $\hatl = (a, b)$, $a \ge b \ge 2$,
   $q \ge a + b - 2$ --- by uniform structural induction.
\item As corollaries: the $(5, 4)$, $(6, 3)$, $(6, 2)$, $(6, 4)$,
   $(5, 5)$, $\ldots$ rows of the ladder.
\end{itemize}

\paragraph{Open.}
\begin{itemize}[topsep=0.3em,itemsep=0.1em]
\item Structural sharpness at $j = \tau(\hatl) + 1$ for general
   two-row $\hatl$.
\item Computational verification at new cases (scripts pending).
\item Multi-row $\hatl$ (length-preserving corner in row $\ge 3$):
   threshold formula needs refinement or new contributor analysis.
\item The dual Frobenius programme (suggested in the May-14
   induced-image refutation note).
\end{itemize}

\end{document}
